%!TEX root = ../Pflichtenheft.tex

\chapter{Software Quality}

\section{Static Analysis}

To ensure a high level of code quality, static code analysis tools were used
for both the client and the server components of the system.

For the client, the TypeScript ecosystem was used. In particular,
\textit{TypeScript type checking} and \textit{ESLint} were applied to detect
potential errors early and to enforce a consistent coding style.
Additionally, the code was automatically formatted to ensure consistent
structure and readability.

For the server, static analysis tools were also used during development.
These tools help identify potential issues such as type inconsistencies,
unused variables, and possible runtime errors before execution.
Regular use of static analysis significantly contributed to maintaining
a high overall code quality throughout the development process.

\section{Extensibility}

The architecture of the system was designed to allow straightforward
extensibility. This is primarily achieved through the clear separation
of responsibilities within a layered architecture.

On the client side, the application follows a structured ViewModel
architecture in which presentation logic is separated from the
underlying business logic. New functionality can therefore often be
added by introducing additional ViewModels or services without
modifying existing components.

Similarly, the server follows a modular architecture. New functionality,
such as additional API endpoints or business logic modules, can be
implemented by adding new modules without affecting existing parts of
the system.

This clear architectural structure ensures that the system can be
extended with additional features in the future without requiring
major refactoring.

\section{Testability}

The system provides a high degree of testability due to the clear
separation between business logic and user interface components.

For the client, extensive automated unit tests were implemented using
the testing framework \textit{Vitest}. In total, 194 tests were executed,
all of which passed successfully. The achieved test coverage exceeds
95\% for statements, lines, and functions.

For the server, automated tests were implemented using the Python
testing framework \textit{pytest}. The test suite contains 142 tests,
all of which passed successfully during execution. The tests are
executed within a Python virtual environment and focus primarily on
the business and service layers of the application.

The server tests cover components such as:

\begin{itemize}
    \item API routes and service endpoints
    \item inference manager
    \item system manager
    \item user information manager
    \item prompt manager
    \item persistence repositories
\end{itemize}

A central test bootstrap was implemented in
\texttt{tests/conftest.py}. This configuration provides default
environment settings and lightweight dependency stubs, enabling the
test suite to run without requiring external services.

The current server test coverage is approximately 63\%. Core components
such as the service APIs, inference logic, system management, and
persistence layer are well covered by the tests. Lower coverage values
mainly originate from modules that depend on external integrations,
such as vector database initialization, graph generation internals,
OpenAI client transport logic, and validator subprocess paths.

During the development of the test suite, several improvements were
also made to the server implementation. These include fixing the
generation of JWT timestamps to use UTC-aware values and correcting
the initialization path for administrator profile pictures.

External dependencies such as network requests or browser APIs are
replaced with mock objects during testing, which allows the tests to
run independently of external services.

Some parts of the user interface, especially complex graphical
components, are not fully covered by automated tests, since they are
difficult to verify meaningfully with unit tests.

\section{Code Readability}

To improve code readability, particular attention was given to clear
structure and consistent naming of classes, methods, and variables.

Additional comments and documentation were added for central
components and complex logic areas. This improves the maintainability
of the code and makes it easier for future developers to understand
and extend the system.

By organizing the system into well-defined modules with clear
interfaces, the code remains understandable and manageable even as
the project grows.

\subsection{Measures to Improve Code Quality}

Several measures were implemented during development to further improve
the overall code quality.

One important step was the introduction of centralized structures for
commonly used configurations and error translations. This prevents
duplicate strings and inconsistent configuration definitions
throughout the codebase.

Furthermore, existing tests were extended and partially restructured
to increase test coverage and improve test stability. In particular,
mock objects were introduced for browser functionality and external
services so that tests can be executed independently from the runtime
environment.

These measures contribute significantly to the maintainability,
stability, and long-term extensibility of the system.